\subsection{Glosario}

\begin{itemize}
    \item \textbf{Alineación Curricular}: Grado de correspondencia, congruencia y consistencia entre los componentes de un programa educativo (tales como perfiles de egreso, planes de estudio y descriptores de asignaturas) con respecto a directrices internacionales y los requerimientos del entorno laboral.
    
    \item \textbf{Clustering (Agrupamiento)}: Técnica de aprendizaje no supervisado que consiste en agrupar un conjunto de objetos de tal manera que los miembros del mismo grupo (\textit{clúster}) sean más similares entre sí que con los de otros grupos, basándose en alguna métrica de distancia o similitud geométrica.
    
    \item \textbf{Coeficiente de Silueta (\textit{Silhouette Coefficient})}: Métrica de validación interna de algoritmos de \textit{clustering} que evalúa qué tan bien agrupado está un elemento. Mide la cercanía de un objeto a su propio clúster (cohesión) frente a su cercanía a otros clústeres (separación), variando entre -1 y 1.
    
    \item \textbf{Corpus}: Conjunto estructurado y homogéneo de textos escritos, recopilado de manera sistemática para ser utilizado como base de análisis lingüístico o procesamiento computacional.
    
    \item \textbf{DBSCAN (\textit{Density-Based Spatial Clustering of Applications with Noise})}: Algoritmo de agrupamiento no supervisado basado en densidad que agrupa puntos cercanos y etiqueta como valores atípicos (\textit{outliers}) aquellos que se encuentran aislados en regiones de baja densidad.
    
    \item \textbf{Embeddings (Incrustaciones Vectoriales)}: Representaciones matemáticas de textos (palabras, frases o documentos completos) en forma de vectores densos en un espacio continuo de baja dimensionalidad, donde la cercanía geométrica de los vectores representa similitud semántica.
    
    \item \textbf{Índice de Davies-Bouldin}: Métrica de evaluación intrínseca de algoritmos de agrupamiento que calcula la similitud promedio entre cada clúster y su clúster más similar. Valores más bajos indican agrupamientos con clústeres más compactos y mejor separados.
    
    \item \textbf{K-Means}: Algoritmo clásico de agrupamiento no supervisado que particiona un conjunto de datos en $K$ clústeres distintos y disjuntos, asignando cada elemento al clúster cuyo centroide se encuentra a menor distancia geométrica (generalmente distancia euclidiana).
    
    \item \textbf{Lematización}: Proceso de preprocesamiento de lenguaje natural que consiste en reducir las palabras de un texto a su forma básica o raíz morfológica canonizada (lema), como infinitivos para verbos y masculino singular para adjetivos y sustantivos.
    
    \item \textbf{Procesamiento de Lenguaje Natural (PLN / NLP)}: Subcampo de la inteligencia artificial y la lingüística computacional enfocado en el diseño e implementación de algoritmos que permiten a las computadoras procesar, entender y analizar textos en lenguaje humano.
    
    \item \textbf{Sentence-Transformers (SBERT)}: Framework de aprendizaje profundo basado en redes neuronales siamesas, utilizado para adaptar modelos de lenguaje tipo BERT para generar embeddings de oraciones y textos completos de alta calidad semántica.
    
    \item \textbf{Similitud del Coseno}: Métrica matemática utilizada para cuantificar la similitud entre dos vectores en un espacio multidimensional, midiendo el coseno del ángulo formado entre ellos, independientemente de su magnitud.
    
    \item \textbf{Stop Words (Palabras Vacías)}: Términos que carecen de significado informativo relevante por sí mismos (como artículos, preposiciones y conjunciones, o en el contexto académico, términos como ``asignatura'', ``horas'' o ``carrera''), que son filtrados durante el preprocesamiento de texto.
    
    \item \textbf{TF-IDF (\textit{Term Frequency - Inverse Document Frequency})}: Técnica estadística clásica de representación de textos que evalúa la importancia de una palabra en un documento en relación con un corpus completo, ponderándola según su frecuencia local y su rareza global.
    
    \item \textbf{t-SNE (\textit{t-Distributed Stochastic Neighbor Embedding})}: Algoritmo no lineal de reducción de dimensionalidad utilizado principalmente para la visualización de datos de alta dimensionalidad en espacios de dos o tres dimensiones.
    
    \item \textbf{UMAP (\textit{Uniform Manifold Approximation and Projection})}: Algoritmo moderno de reducción de dimensionalidad no lineal fundamentado en la geometría de variedades, eficiente computacionalmente y robusto para preservar la estructura local y global de los datos.
\end{itemize}